\chapter{Introduction}

The notion of defects in quantum field theories (QFTs) has amassed an extraordinary amount of attention and importance in physics and, perhaps surprisingly, also in mathematics. Let $M$ be a smooth $n$-dimensional manifold on which a given QFT $\tau$ lives. A defect of dimension $d$ can be heuristically articulated as the data of a submanifold $D\subset M$ of codimension $(n-d)$, together with some discontinuity of $\tau$ along $D$ which produces another quantum field theory. Examples of these phenomena are given by Wilson and ’t Hooft line defects in Yang--Mills theory. Giving a definition and a classification of these general defects is the subject of extensive ongoing work (see \cite{CEG23,CZR25}).

QFTs which lack dependence on any geometric parameter are dubbed \emph{topological QFTs} (TQFTs). In particular, defects whose interaction with the bulk theory is independent of the isotopy class of their embedding are called \emph{topological defects}. One can organize topological QFTs with topological defects into a single package, called a defect topological quantum field theory, as they are defined in the work of Carqueville et al. — for an overview of recent developments see \cite{CZR25} and for a thorough introduction see \cite{CRS19}. These are symmetric monoidal functors out of the 1-category of (oriented) stratified cobordisms with defect data into the category of complex vector spaces
\begin{equation}
\cZ \maps \mathsf{Bord}_{d}(\DD) \longra \mathsf{Vect}_{\CC},
\end{equation}
where $\DD=(D_d,\ldots,D_0)$ is the collection of defect labels — when omitting the defect labels $\DD$, we will call these \emph{stratified} TQFTs.

Extending the categories of defect bordisms upwards — by adding stratified diffeomorphisms and stratified isotopies between them, and higher stratified isotopies between these, etc. — we can contemplate derived versions (or $\oo$ versions) of these defect TQFTs; whereas by extending them downwards — letting stratified bordisms be cut into stratified cornered pieces of arbitrary depth — we can contemplate extended and more computable versions of these derived defect TQFTs. In this thesis we will construct the $(\oo,1)$-category of stratified bordisms whose objects are closed $(d-1)$-dimensional topological manifolds endowed with a conically smooth stratified atlas; the 1-morphisms are $d$-dimensional bordisms endowed also with a conically smooth atlas; the higher morphisms encode conically smooth isomorphisms and conically smooth isotopies among them. We consider also versions with defect labels. In particular, we sketch a way of interpreting the defect labels $\DD$ defined in \cite[Definition~2.4]{CRS19} as part of the tangential data associated to a conically smooth stratified space. Lastly, we provide a sketch of a definition of the $(\oo,d)$-category of fully-extended stratified bordisms: objects are points; $k$-morphisms are stratified $k$-bordisms for $k\leq d$ and stratified isotopies for $k > d$.     


Despite the growing interest in defect TQFTs, especially derived and extended ones, the literature on such $(\oo,d)$-categories of stratified bordisms (or defect bordisms) is scant still, and mostly limited to specific ranges of codimension for the stratification and to low dimensions. Our construction is based on the following three choices:
\begin{enumerate}
\item We use a different model for higher categories, which we call complete ``interval'' Segal spaces (\Cref{intro:: thm: interval models}).
\item We make use of the theory of conically smooth stratified spaces and their tangential structures — based on the work of Ayala--Francis--Tanaka--Rozenblyum (\Cref{intro:: thm: stratified bordisms}).
\item We limit our construction to stratified bordisms whose underlying space is a topological manifold.
\end{enumerate}
These choices mean that our construction, while close in spirit to Lurie's in \cite[Definition~4.11]{Lur09}, explores a more geometrically-biased variant of Segal spaces, enjoys a more solid differential-topological foundation coming from the comprehensive works of Ayala--Francis--Tanaka \cite{AFT1} and Ayala--Francis--Rozenblyum \cite{AFR1}\cite{AFR2}, and is less general insofar as it concerns only topological spaces whose underlying space is locally Euclidean and of fixed dimension. This last point, we believe, is not a true limitation of our construction and we intend for future work to address it — concretely, the cornerstone of our approach is a version of the stratified Ehresmann fibration theorem, which does not apply only to conically smooth stratified spaces whose underlying topological space is locally Euclidean.

\section*{Main results}

We now state the three principal results of the thesis. The key idea is to
replace the simplicial indexing category $\bDelta^{\op}$ by a geometric
category $\Int$. Informally, an object of $\Int$ is a finite union of open
subintervals of $[0,1]$ containing both endpoints, but not the whole interval;
morphisms are inclusions. Its complement is a non-empty ordered collection
of closed intervals, which specify regions over which a bordism can be cut
and glued. The functor $\gamma\colon\Int\to\bDelta^{\op}$ remembers only
the ordered set of these complementary components.

A \emph{pre-interval space} is a functor $\Int\to\sset$, where $\sset$
denotes simplicial sets. An \emph{interval Segal space} sends refinements
preserving the cutting components to weak equivalences and reconstructs
chains of morphisms from their individual pieces by homotopy pullbacks.
Completeness identifies the space of objects with the space of equivalences,
as for Rezk's complete Segal spaces.

\begin{samepage}
\begin{maintheorem}[Interval models for $(\oo,1)$-categories]
\label{intro:: thm: interval models}
Let $\Int$ be the inclusion poset of finite unions of open subintervals of
$[0,1]$ containing $\{0,1\}$ and different from $[0,1]$. The category
$\fun(\Int,\sset)$ of pre-interval spaces admits a left proper, combinatorial, and
simplicial model structure whose fibrant objects are the complete interval
Segal spaces. Moreover, restriction and right Kan extension along
\[
    \gamma\colon \Int\longrightarrow\bDelta^{\op}
\]
where $\gamma$ records the ordered complementary components, induce a Quillen equivalence
\[
    \mathsf{CSSp}
    \ \underset{\rkan_{\gamma}}{\overset{\gamma^{*}}{\rightleftarrows}}\
    \mathsf{CISSp}.
\]
Here $\mathsf{CSSp}$ is Rezk's model category of complete Segal spaces and
$\mathsf{CISSp}$ is the model category just described.
Consequently, complete interval Segal spaces present the
$(\oo,1)$-category of $(\oo,1)$-categories.
\end{maintheorem}
\end{samepage}

The significance of \Cref{intro:: thm: interval models} is that replacing
$\bDelta^{\op}$ by $\Int$ retains geometric cutting data without changing
the homotopy theory of $(\oo,1)$-categories. The precise statement and proof
are given in \Cref{ch2:: thm: complete interval model}.

For the geometric constructions we use the smooth-sheaf framework of
Berwick-Evans, Boavida de Brito, and Pavlov (BBEP) \cite{BBEP24}.
The \emph{smooth moduli sheaf of embedded bordisms} assigns to a smooth
parameter manifold the smooth families of embedded bordisms over it, with
specified cylindrical cutting regions. Its \emph{concordance resolution}
organizes these families into a simplicial space: vertices are individual
bordisms, paths are smooth one-parameter families, and higher simplices
record families with more parameters. This retains deformation data rather
than taking only a set of diffeomorphism classes. Stabilizing the ambient
Euclidean space means allowing arbitrarily many embedding coordinates.

\begin{samepage}
\begin{maintheorem}[The smooth bordism $(\oo,1)$-category]
\label{intro:: thm: smooth bordisms}
For every $d\geq 0$, form $\bord_d\colon\Int\to\sset$ from smooth families
of embedded $d$-dimensional bordisms with cylindrical cutting regions by
the concordance resolution of Berwick-Evans, Boavida de Brito, and Pavlov
\cite{BBEP24} and stabilization of the ambient Euclidean
space. Then $\bord_d$ is a symmetric monoidal interval Segal space, with
monoidal operation given by disjoint union. Its objects are closed smooth
$(d-1)$-manifolds, its morphisms are smooth $d$-bordisms, and its higher
morphisms encode diffeomorphisms and their isotopies. Its homotopy category
$\ho(\bord_d)$, obtained by identifying homotopic morphisms, is equivalent
as a symmetric monoidal category to the ordinary unoriented bordism category
$\mathsf{Cob}_d$, whose morphisms are diffeomorphism classes of bordisms:
\[
    \ho(\bord_d)\simeq\mathsf{Cob}_d.
\]
\end{maintheorem}
\end{samepage}

Thus the interval model supports a geometric construction of the bordism
$(\oo,1)$-category in which smooth families are retained rather than
prematurely divided out by diffeomorphisms. The interval Segal condition,
symmetric monoidal structure, and homotopy-category comparison are proved
in \Cref{ch3:: thm: interval segal bordism cat,ch3:: thm: symmetric monoidal,ch3:: thm: homotopy category bordism}, respectively.

\begin{samepage}
\begin{maintheorem}[The stratified bordism $(\oo,1)$-category]
\label{intro:: thm: stratified bordisms}
For every $d\geq 0$, form $\bordstr_d\colon\Int\to\sset$ from smooth
families of embedded conically smooth stratified topological $d$-bordisms
with cylindrical cutting regions by the concordance resolution of
Berwick-Evans, Boavida de Brito, and Pavlov \cite{BBEP24} and
stabilization of the ambient Euclidean space. Here the underlying spaces
are topological manifolds and the stratifications have conically smooth
local models. Then $\bordstr_d$ is a symmetric monoidal interval Segal space,
with monoidal operation given by disjoint union. Its objects are closed
stratified topological $(d-1)$-manifolds, its morphisms are stratified
$d$-bordisms, and its higher morphisms encode conically smooth isomorphisms
and their isotopies. Its homotopy category is equivalent as a symmetric
monoidal category to $\cS\mathsf{Cob}_d$, the ordinary unoriented stratified
bordism category whose morphisms are conically smooth isomorphism classes
of bordisms:
\[
    \ho(\bordstr_d)\simeq\cS\mathsf{Cob}_d.
\]
\end{maintheorem}
\end{samepage}

The construction and comparison in \Cref{intro:: thm: stratified bordisms}
are proved in \Cref{ch5:: thm: stratbord interval segal,ch5:: thm: symmetric monoidal stratified,ch5:: thm: homotopy category stratified bord}.
They rest on several geometric results established in
Chapter~4. We prove that every proper conically smooth submersion is a
conically smooth fibre bundle, together with the corresponding
boundary-compatible statement needed for bordisms. We also prove that every
finite conically smooth stratified space admits a proper conically smooth
embedding into some Euclidean space and that the stabilized embedding space
\[
    \Strat^{\sf emb}(M,\RR^{\oo})
\]
is weakly contractible. For compact depth-$1$ spaces of dimension $n$, we
obtain the sharper embedding dimension $\RR^{2n+1}$. These results supply the
local triviality and embedding-independence required in the construction of
$\bordstr_d$.

The remainder of the introduction explains the three choices underlying
these results and places them in relation to earlier constructions.

\emph{Why do we use a different model for higher categories?} The theory of $(\oo,d)$-categories is already filled with an abundance of papers and books covering a variety of different models for $(\oo,d)$-categories. Ours is extremely close to that of complete Segal spaces — first articulated in \cite{Rezk01}.\footnote{A thorough introduction to the zoo of models for the homotopy theory of $(\oo,1)$-categories, we highly recommend the comprehensive volume of Julia Bergner \cite{Ber18}} Recall that a complete Segal space is a pre-sheaf 
\begin{equation}
    \cC \maps \bDelta^{\op} \longra \sset
\end{equation}
which satisfies the Segal condition 
\begin{equation}
    \cC([n]) \,\, \simeq \,\, \cC([1]) \times_{\cC([0])} \cC([1]) \times_{\cC([0])} \cdots \times_{\cC([0])} \cC([1]).
\end{equation}
Since the collection of all objects $\cC([0])$ forms not a set but a space, the \emph{completeness condition} merely ensures that this space agrees with the space of equivalences of $\cC$, so category theory and homotopy theory work in tandem. We introduce a variant of the simplex category $\bDelta$ by replacing finite linearly ordered sets with something slightly more geometric: open intervals in the closed interval $[0,1]$. We introduce a category $\Int$, whose objects are finite disjoint unions of open subintervals containing both endpoints, but different from $[0,1]$; morphisms are inclusions. This admits a functor to $\bDelta^{\op}$ which can be sketched as follows:
\begin{equation}\label{intro:: eq: interval_object}
    I \,\, = \,\, \tikzfig{figures/Ch2/INT1}{\begin{tikzpicture}
	\begin{pgfonlayer}{nodelayer}
		\node [style=none] (0) at (-7, 0) {};
		\node [style=none] (1) at (8, 0) {};
		\node [style=none] (2) at (-5, 0) {};
		\node [style=none] (3) at (-3, 0) {};
		\node [style=none] (4) at (-2.5, 0) {};
		\node [style=none] (5) at (-1.5, 0) {};
		\node [style=none] (6) at (2.5, 0) {};
		\node [style=none] (7) at (6.75, 0) {};
		\node [style=endpoint] (8) at (-7, 0) {};
		\node [style=endpoint] (9) at (8, 0) {};
		\node [style=none] (10) at (-7, 0.75) {$0$};
		\node [style=none] (11) at (8, 0.75) {$1$};
	\end{pgfonlayer}
	\begin{pgfonlayer}{edgelayer}
		\draw (0.center) to (1.center);
		\draw [style={open_int}] (0.center) to (2.center);
		\draw [style={open_int}] (3.center) to (4.center);
		\draw [style={open_int}] (5.center) to (6.center);
		\draw [style={open_int}] (7.center) to (1.center);
	\end{pgfonlayer}
\end{tikzpicture}
}
\end{equation}
We prove that functors from $\Int$ to $\sset$, mapping inclusions which preserve connected components to weak equivalences, form a model category that is Quillen equivalent to a projective/injective model category of functors from $\bDelta^{\op}$ to $\sset$ — where in both cases the weak equivalences are defined levelwise. This is \emph{almost} equivalent to stating that $\bDelta^{\op}$ is an $\oo$-localization of $\Int$ relative to component-preserving inclusions. From this Quillen equivalence we articulate an adapted Segal condition which we call interval Segal condition, and prove that Segal spaces are Quillen equivalent to interval Segal spaces.

The reason we develop this model for $(\oo,1)$-categories, instead of using complete Segal spaces, is twofold:
\begin{enumerate}
    \item It is intrinsically geometric in nature, which better aligns with our intent of constructing an $(\oo,1)$-category of cobordisms. In past constructions, this interval structure has always played an adjacent role rather than an intrinsic one. It is the author's belief that interval Segal spaces offer a more suitable canvas for building $(\oo,1)$-categories of a geometric nature.
    \item It generalizes well for higher categories, i.e.,\ $(\oo,d)$-categories. As we expand on in the Further Outlook chapter, at the end of the thesis, $\Int$ admits a natural extension to cubes which can just as well parametrize $\btheta_d$ or $\bDelta^{\times d}$ — serving as appropriate shapes to model $(\oo,d)$-categories or cubical versions thereof.
\end{enumerate}

A last --- yet paramount — technical reason for our shift in model, has to do with how well it blends with the spatial structure, i.e.,\ the $\oo$ part of the definition of bordism categories. It has been diagnosed, first in the work of Violeta Marques and Gustavo Granja, that the most state-of-the-art definition of $(\oo,d)$-category of fully-extended cobordisms, given by Claudia Scheimbauer and Damien Calaque \cite{CS19}, suffers from two subtle technical flaws: 

\begin{enumerate}
    \item It assumes that the smooth singular complex of a smooth manifold with corners is a Kan complex — which is proved not to be case in the work of Dan Christensen and Enxin Wu \cite{CW}. See \cite[Section~4.1]{CS19}.
    \item It states that the smooth singular simplices of the topological space of submanifolds satisfy a local triviality condition which would ensure that they form a Kan complex, but they do not satisfy this condition. See \cite[Lemma~5.7]{CS19}.
\end{enumerate}

To provide some additional context for the importance of this issue, we will first discuss how modern constructions of higher categories of bordisms are usually set up. We are used to formulating bordism categories from the \emph{bottom up}: in other words, we define the bordism category as a $k$-category (or $(\oo,k)$-category) whose $k$-morphisms are $(d-k)$-dimensional manifolds with corners and whose composition is given by gluing along the ``right'' boundary components. While this interpretation is retained in the works of Lurie and Calaque--Scheimbauer, their approach can be described as \emph{top down}. As is well documented in the literature, one of the most contrived aspects of constructing higher categories of bordisms is defining composition, since it comes with a compounding pile of choices whose compatibility is easy to check, but very hard to write down. We will first go over a sketch of the definition of Lurie, then we will explain how Calaque--Scheimbauer (CS) address the problems in Lurie's construction and expand on how ours fixes the issues with CS indicated above.

Sticking to the $(\oo,1)$-case, for simplicity, one might describe a point in Lurie's space of $k$-composable chains of 1-morphisms as a choice of a finite-dimensional real vector space $V$ and an embedding $M \hookra V \times \RR$, where $M$ is a $d$-dimensional smooth manifold together with a set $\{t_0\leq \cdots \leq t_k\} \subset \RR$ of $k$ points whose restricted projection  $\pi_{|M}\maps M \fib \RR$ is proper and such that
\begin{equation}
    \mathsf{Crit}(\pi_{|M}) \cap \{t_0\leq \cdots \leq t_k\} = \emptyset.
\end{equation}
Here $\mathsf{Crit}(\pi_{|M})\subset\RR$ denotes the set of critical
\emph{values} of the restricted projection. We depict one such example for $k=2$
\begin{equation}
    \scalebox{0.8}{\tikzfig{figures/Ch3/BOR6}{\begin{tikzpicture}
	\begin{pgfonlayer}{nodelayer}
		\node [style=none] (0) at (-9.5, 0) {};
		\node [style=none] (1) at (-4.5, 3) {};
		\node [style=none] (2) at (-4.5, -3) {};
		\node [style=none] (3) at (9.5, 0) {};
		\node [style=none] (4) at (-4.5, 3) {};
		\node [style=none] (5) at (-4.5, -3) {};
		\node [style=none] (6) at (-7, 0.25) {};
		\node [style=none] (7) at (-2, 0.25) {};
		\node [style=none] (8) at (-6.75, -0.35) {};
		\node [style=none] (9) at (-2.5, -0.425) {};
		\node [style=none] (10) at (1.5, 0.25) {};
		\node [style=none] (11) at (6.5, 0.25) {};
		\node [style=none] (12) at (1.75, -0.35) {};
		\node [style=none] (13) at (6, -0.425) {};
		\node [style=none] (14) at (-12, 7) {};
		\node [style=none] (15) at (-12, -6.5) {};
		\node [style=none] (16) at (12, -6) {};
		\node [style=none] (17) at (-12.5, -6) {};
		\node [style=none] (18) at (-13, 7) {$V$};
		\node [style=none] (19) at (12, -6.75) {};
		\node [style=none] (20) at (-6.75, 6.5) {};
		\node [style=none] (21) at (-3.5, 9) {};
		\node [style=none] (22) at (-6.75, -7.5) {};
		\node [style=none] (23) at (-3.5, -5) {};
		\node [style=none] (32) at (-4.75, -6) {};
		\node [style=none] (33) at (-2, 6.5) {};
		\node [style=none] (34) at (1.25, 9) {};
		\node [style=none] (35) at (-2, -7.5) {};
		\node [style=none] (36) at (1.25, -5) {};
		\node [style=none] (37) at (2.25, 6.5) {};
		\node [style=none] (38) at (5.5, 9) {};
		\node [style=none] (39) at (2.25, -7.5) {};
		\node [style=none] (40) at (5.5, -5) {};
		\node [style=none] (41) at (0, -6) {};
		\node [style=none] (42) at (4.25, -6) {};
		\node [style=none] (43) at (-4.75, -7) {$t_0$};
		\node [style=none] (44) at (0, -7) {$t_1$};
		\node [style=none] (45) at (4.25, -7) {$t_2$};
		\node [style=none] (46) at (-4.75, -3) {};
		\node [style=none] (47) at (-4.75, -0.85) {};
		\node [style=none] (48) at (-4.75, 0.55) {};
		\node [style=none] (49) at (-4.75, 3) {};
		\node [style=none] (50) at (4.25, -3.65) {};
		\node [style=none] (51) at (4.25, -0.85) {};
		\node [style=none] (52) at (4.25, 0.55) {};
		\node [style=none] (53) at (4.325, 3.625) {};
		\node [style=none] (54) at (0, -3.325) {};
		\node [style=none] (55) at (-0.05, 3.325) {};
	\end{pgfonlayer}
	\begin{pgfonlayer}{edgelayer}
		\draw [style=geometric, in=90, out=-180] (1.center) to (0.center);
		\draw [style=geometric, in=180, out=-90] (0.center) to (2.center);
		\draw [style=geometric, in=90, out=0] (4.center) to (3.center);
		\draw [style=geometric, in=0, out=-90] (3.center) to (5.center);
		\draw [style=geometric, in=270, out=-105, looseness=0.75] (6.center) to (7.center);
		\draw [style=geometric, in=90, out=90, looseness=0.75] (8.center) to (9.center);
		\draw [style=geometric, in=270, out=-105, looseness=0.75] (10.center) to (11.center);
		\draw [style=geometric, in=90, out=90, looseness=0.75] (12.center) to (13.center);
		\draw [style=dir] (15.center) to (14.center);
		\draw [style=dir] (17.center) to (16.center);
		\draw [style=geometric,line width=1.05pt] (-4.75, -6.20) to (-4.75, -5.80);
		\draw [style=geometric,line width=1.05pt] (0, -6.20) to (0, -5.80);
		\draw [style=geometric,line width=1.05pt] (4.25, -6.20) to (4.25, -5.80);
		\draw [style=dir,line width=1.05pt] (-4.75, -4.45) to (-4.75, -5.62);
		\draw [style=dir,line width=1.05pt] (0, -4.45) to (0, -5.62);
		\draw [style=dir,line width=1.05pt] (4.25, -4.45) to (4.25, -5.62);
		\draw [in=180, out=-180, looseness=0.75] (49.center) to (48.center);
		\draw [in=180, out=-180, looseness=0.75] (47.center) to (46.center);
		\draw [in=180, out=-180, looseness=0.50] (53.center) to (52.center);
		\draw [in=180, out=-180, looseness=0.50] (51.center) to (50.center);
		\draw [in=180, out=-180, looseness=0.25] (55.center) to (54.center);
	\end{pgfonlayer}
\end{tikzpicture}
}}
\end{equation} 
\medskip

This set of embeddings can be topologized appropriately and a valid model for the correct $(\oo,1)$-category of bordisms might be extracted. Despite this, the Segal condition proves hard to check, and it does so for reasons which are unsurprising at this stage: cutting manifolds across a single submanifold, or gluing them if you see it as the reverse operation, is easier to do if you first specify collars. Failure to do so begets issues of coherence and pseudo-isotopy — which, although not intractable, increase the burden of working with such models. This issue can be ameliorated by placing the legwork elsewhere, and this is at the core of the variant introduced in the work of Calaque--Scheimbauer. 

To build collars into the definition, Calaque--Scheimbauer introduce an auxiliary $(\oo,1)$-category of closed intervals $\cI$. Calaque--Scheimbauer denote this category by $\mathrm{Int}$; we use $\cI$ to distinguish it from our interval category $\Int$. Assuming, for the sake of simplicity, that there were no issues with this definition, then as a complete Segal space its $k$-levels are defined as $(k+1)$-tuples 
\begin{equation}
    \cI_k = \{([a_0,b_0]\leq \cdots \leq [a_k,b_k])\,:\, [a_i,b_i]\subset \RR \}
\end{equation}
topologized via the standard Euclidean topology on $\RR^{2k+2}$,  and this is where the technical issue appears — in order for these spaces to organize into a Segal space, their smooth singular complex must be a Kan complex, and for $k > 1$ these are not; as $\cI_{k}$ is a smooth manifold with corners. Working with these base intervals, it improves on Lurie's definition by pairing the embedding $M \hookra V\times \RR$ with a choice of a $k$-simplex $(I_0\leq \cdots \leq I_k)\in \cI_k$ such that $\mathsf{Crit}_{|\pi}$ is instead disjoint from $I_0 \cup \cdots \cup I_k$. Visually, we can depict a $k$-simplex as follows 
\begin{equation}
    \scalebox{0.8}{\tikzfig{figures/Ch3/BOR7}{\begin{tikzpicture}
	\begin{pgfonlayer}{nodelayer}
		\node [style=none] (0) at (-9.5, 0) {};
		\node [style=none] (1) at (-4.5, 3) {};
		\node [style=none] (2) at (-4.5, -3) {};
		\node [style=none] (3) at (9.5, 0) {};
		\node [style=none] (4) at (-4.5, 3) {};
		\node [style=none] (5) at (-4.5, -3) {};
		\node [style=none] (6) at (-7, 0.25) {};
		\node [style=none] (7) at (-2, 0.25) {};
		\node [style=none] (8) at (-6.75, -0.35) {};
		\node [style=none] (9) at (-2.5, -0.425) {};
		\node [style=none] (10) at (1.5, 0.25) {};
		\node [style=none] (11) at (6.5, 0.25) {};
		\node [style=none] (12) at (1.75, -0.35) {};
		\node [style=none] (13) at (6, -0.425) {};
		\node [style=none] (14) at (-12, 7) {};
		\node [style=none] (15) at (-12, -6.5) {};
		\node [style=none] (16) at (12, -6) {};
		\node [style=none] (17) at (-12.5, -6) {};
		\node [style=none] (18) at (-13, 7) {$V$};
		\node [style=none] (19) at (12, -6.75) {};
		\node [style=none] (32) at (-4.75, -6) {};
		\node [style=none] (41) at (0, -6) {};
		\node [style=none] (42) at (4.25, -6) {};
		\node [style=none] (46) at (-4.75, -3) {};
		\node [style=none] (47) at (-4.75, -0.85) {};
		\node [style=none] (48) at (-4.75, 0.55) {};
		\node [style=none] (49) at (-4.75, 3) {};
		\node [style=none] (50) at (4.25, -3.65) {};
		\node [style=none] (51) at (4.25, -0.85) {};
		\node [style=none] (52) at (4.25, 0.55) {};
		\node [style=none] (53) at (4.325, 3.625) {};
		\node [style=none] (54) at (0, -3.325) {};
		\node [style=none] (55) at (-0.05, 3.325) {};
		\node [style=none] (56) at (-5.75, -6) {};
		\node [style=none] (57) at (-4, -6) {};
		\node [style=none] (58) at (-1, -6) {};
		\node [style=none] (59) at (1.25, -6) {};
		\node [style=none] (60) at (2.75, -6) {};
		\node [style=none] (61) at (5, -6) {};
		\node [style=none] (62) at (-5.75, 0.475) {};
		\node [style=none] (63) at (-4, 0.525) {};
		\node [style=none] (64) at (-5.875, 2.9) {};
		\node [style=none] (65) at (-3.75, 3.025) {};
		\node [style=none] (66) at (-5.85, -2.925) {};
		\node [style=none] (67) at (-4.1, -3) {};
		\node [style=none] (68) at (-5.925, -0.7) {};
		\node [style=none] (69) at (-3.85, -0.825) {};
		\node [style=none] (70) at (-1.1, -3.225) {};
		\node [style=none] (71) at (1.15, -3.475) {};
		\node [style=none] (72) at (-1, 3.225) {};
		\node [style=none] (73) at (1.075, 3.425) {};
		\node [style=none] (74) at (2.65, -3.6) {};
		\node [style=none] (75) at (5.025, -3.65) {};
		\node [style=none] (76) at (2.75, -0.75) {};
		\node [style=none] (77) at (5.125, -0.75) {};
		\node [style=none] (78) at (2.8, 0.475) {};
		\node [style=none] (79) at (5.15, 0.4) {};
		\node [style=none] (80) at (2.825, 3.6) {};
		\node [style=none] (81) at (5.15, 3.625) {};
		\node [style=none] (82) at (-5.75, -6.75) {$a_0$};
		\node [style=none] (83) at (-4, -6.75) {$b_0$};
		\node [style=none] (84) at (-5.75, -5.75) {};
		\node [style=none] (85) at (-5.75, -6.25) {};
		\node [style=none] (86) at (-4, -5.75) {};
		\node [style=none] (87) at (-4, -6.25) {};
		\node [style=none] (88) at (-1, -5.75) {};
		\node [style=none] (89) at (-1, -6.25) {};
		\node [style=none] (90) at (1.25, -5.75) {};
		\node [style=none] (91) at (1.25, -6.25) {};
		\node [style=none] (92) at (2.75, -5.75) {};
		\node [style=none] (93) at (2.75, -6.25) {};
		\node [style=none] (94) at (5, -5.75) {};
		\node [style=none] (95) at (5, -6.25) {};
		\node [style=none] (96) at (-1, -6.75) {$a_1$};
		\node [style=none] (97) at (1.25, -6.75) {$b_1$};
		\node [style=none] (98) at (2.75, -6.75) {$a_2$};
		\node [style=none] (99) at (5, -6.75) {$b_2$};
	\end{pgfonlayer}
	\begin{pgfonlayer}{edgelayer}
		\draw [style=geometric, in=90, out=-180] (1.center) to (0.center);
		\draw [style=geometric, in=180, out=-90] (0.center) to (2.center);
		\draw [style=geometric, in=90, out=0] (4.center) to (3.center);
		\draw [style=geometric, in=0, out=-90] (3.center) to (5.center);
		\draw [style=geometric, in=270, out=-105, looseness=0.75] (6.center) to (7.center);
		\draw [style=geometric, in=90, out=90, looseness=0.75] (8.center) to (9.center);
		\draw [style=geometric, in=270, out=-105, looseness=0.75] (10.center) to (11.center);
		\draw [style=geometric, in=90, out=90, looseness=0.75] (12.center) to (13.center);
		\draw [style=dir] (15.center) to (14.center);
		\draw [style=dir] (17.center) to (16.center);
		\draw [in=180, out=-180, looseness=0.50] (49.center) to (48.center);
		\draw [in=180, out=-180, looseness=0.75] (47.center) to (46.center);
		\draw [in=180, out=-180, looseness=0.50] (53.center) to (52.center);
		\draw [in=180, out=-180, looseness=0.50] (51.center) to (50.center);
		\draw [in=180, out=-180, looseness=0.25] (55.center) to (54.center);
		\draw [style={open_int}] (56.center) to (57.center);
		\draw [style={open_int}] (59.center) to (58.center);
		\draw [style={open_int}] (60.center) to (61.center);
		\draw [style=red shade] (65.center)
			 to [in=30, out=-180, looseness=0.25] (64.center)
			 to [in=-165, out=-165, looseness=0.50] (62.center)
			 to [in=165, out=15, looseness=0.50] (63.center)
			 to [in=180, out=180, looseness=0.50] cycle;
		\draw [style=red shade] (69.center)
			 to [in=-30, out=-180, looseness=0.50] (68.center)
			 to [in=-165, out=-165, looseness=0.50] (66.center)
			 to [in=180, out=0, looseness=0.50] (67.center)
			 to [in=-180, out=-180, looseness=0.50] cycle;
		\draw [style=red shade] (71.center)
			 to [in=180, out=180, looseness=0.25] (73.center)
			 to [in=0, out=-180, looseness=0.25] (72.center)
			 to [in=-180, out=-180, looseness=0.25] (70.center)
			 to [in=165, out=0, looseness=0.50] cycle;
		\draw [style=red shade] (75.center)
			 to [in=-165, out=180, looseness=0.50] (77.center)
			 to [in=-15, out=-165, looseness=0.50] (76.center)
			 to [in=-165, out=-165, looseness=0.50] (74.center)
			 to [in=180, out=0, looseness=0.50] cycle;
		\draw [style=red shade] (80.center)
			 to [in=-165, out=-165, looseness=0.50] (78.center)
			 to [in=165, out=15, looseness=0.50] (79.center)
			 to [in=180, out=-180, looseness=0.50] (81.center)
			 to [in=0, out=165, looseness=0.25] cycle;
		\draw [style=red shade, in=135, out=45, loop] (56.center) to ();
		\draw [style=red shade, in=135, out=45, loop] (56.center) to ();
		\draw [style=red shade, in=135, out=45, loop] (56.center) to ();
		\draw [style=red shade, in=135, out=45, loop] (85.center) to ();
		\draw [style={open_int}] (84.center) to (85.center);
		\draw [style=red shade, in=135, out=45, loop] (87.center) to ();
		\draw [style={open_int}] (86.center) to (87.center);
		\draw [style=red shade, in=135, out=45, loop] (89.center) to ();
		\draw [style={open_int}] (88.center) to (89.center);
		\draw [style=red shade, in=135, out=45, loop] (91.center) to ();
		\draw [style={open_int}] (90.center) to (91.center);
		\draw [style=red shade, in=135, out=45, loop] (93.center) to ();
		\draw [style={open_int}] (92.center) to (93.center);
		\draw [style=red shade, in=135, out=45, loop] (95.center) to ();
		\draw [style={open_int}] (94.center) to (95.center);
	\end{pgfonlayer}
\end{tikzpicture}
}}
\end{equation}
\medskip

\nid where the red shaded regions in $M$ are manifestly the pre-images of $\pi_{|M}$ over the intervals $I_i$, and correspond to the collars along which the bordism decomposes. This definition has a slightly larger amount of overhead than Lurie's, but it offers a clearer way of checking the Segal condition and it fixes a failure of essential constancy in Lurie's model. In particular, it is claimed that the spatial $l$-simplices are obtained as the smooth singular complex of a moduli space of embedded submanifolds crossed with $\cI_{k}$. There are two issues here: first, it is not clear to us that what smooth structure is being assumed for the moduli space of embedded submanifolds, calling into question whether the resulting simplicial set is a Kan complex; second, it is stated that these $l$-simplices can be described as the those satisfying conditions (1)--(3) in \cite[Definition~5.6]{CS19}, but it is not clear how condition 2b. would result from taking smooth singular simplices, and we argue that if condition 2b. is assumed then the Kan condition must fail, as one could find examples of horn conditions which do not admit fillings. This issue lies at the core of our choice to change both the model for $(\oo,1)$-category as well as our approach to the spatial structure.



Our definition tries to reduce the extra machinery in Calaque--Scheimbauer's interval $(\oo,1)$-category $\cI$ by, once again, placing it elsewhere. We use interval Segal spaces and we no longer talk about $k$-simplices as direct encoders of chains of composable morphisms, instead we use a disjoint union of open intervals $I \subset [0,1]$ containing the endpoints. An $I$-level element will be, similarly to the above constructions, an embedding $M \hookra V\times [0,1]$, which we depict as
\begin{equation}
    \scalebox{0.8}{\tikzfig{figures/Ch3/BOR8}{\begin{tikzpicture}
	\begin{pgfonlayer}{nodelayer}
		\node [style=none] (1) at (-4.5, 3) {};
		\node [style=none] (2) at (-4.5, -3) {};
		\node [style=none] (3) at (9.5, 0) {};
		\node [style=none] (4) at (-4.5, 3) {};
		\node [style=none] (5) at (-4.5, -3) {};
		\node [style=none] (7) at (-2, 0.25) {};
		\node [style=none] (10) at (1.5, 0.25) {};
		\node [style=none] (11) at (6.5, 0.25) {};
		\node [style=none] (12) at (1.75, -0.35) {};
		\node [style=none] (13) at (6, -0.425) {};
		\node [style=none] (14) at (-12, 7) {};
		\node [style=none] (15) at (-12, -6.5) {};
		\node [style=none] (16) at (13.75, -6) {};
		\node [style=none] (17) at (-12.5, -6) {};
		\node [style=none] (18) at (-13, 7) {$V$};
		\node [style=none] (41) at (0, -6) {};
		\node [style=none] (42) at (4.25, -6) {};
		\node [style=none] (46) at (-4.75, -3) {};
		\node [style=none] (47) at (-4.75, -0.85) {};
		\node [style=none] (48) at (-4.75, 0.55) {};
		\node [style=none] (49) at (-4.75, 3) {};
		\node [style=none] (50) at (4.25, -3.65) {};
		\node [style=none] (51) at (4.25, -0.85) {};
		\node [style=none] (52) at (4.25, 0.55) {};
		\node [style=none] (53) at (4.325, 3.625) {};
		\node [style=none] (55) at (-0.05, 3.325) {};
		\node [style=none] (57) at (-4, -6) {};
		\node [style=none] (58) at (-1, -6) {};
		\node [style=none] (59) at (1.25, -6) {};
		\node [style=none] (60) at (2.75, -6) {};
		\node [style=none] (61) at (5, -6) {};
		\node [style=none] (81) at (5, 3.625) {};
		\node [style=none] (82) at (-11, -6) {};
		\node [style=none] (83) at (-12, -6) {};
		\node [style=none] (84) at (-10.25, -6) {};
		\node [style=none] (85) at (-6, -6) {};
		\node [style=none] (86) at (10, -6) {};
		\node [style=none] (87) at (11.25, -6) {};
		\node [style=none] (88) at (13, -6) {};
		\node [style=none] (90) at (-5, 0.5) {};
		\node [style=none] (92) at (-7.25, 0.25) {};
		\node [style=none] (93) at (-9.5, -10) {};
		\node [style=none] (100) at (-9.5, 0) {};
		\node [style=none] (101) at (-6.95, -0.225) {};
		\node [style=none] (102) at (-8.175, -2.225) {};
		\node [style=none] (103) at (-7.95, 2.325) {};
		\node [style=none] (104) at (-6, 2.925) {};
		\node [style=none] (105) at (-6, 0.4) {};
		\node [style=none] (106) at (-6, -0.725) {};
		\node [style=none] (107) at (-6, -2.875) {};
		\node [style=none] (108) at (9.5, 0) {};
		\node [style=none] (109) at (6.025, -0.4) {};
		\node [style=none] (110) at (7.975, 2.95) {};
		\node [style=none] (111) at (7.825, -3.075) {};
		\node [style=none] (112) at (5, -3.65) {};
		\node [style=none] (113) at (5, -0.75) {};
		\node [style=none] (114) at (5, 0.425) {};
		\node [style=none] (115) at (5, 3.625) {};
		\node [style=none] (116) at (-2.5, -0.4) {};
		\node [style=none] (118) at (-1, 3.25) {};
		\node [style=none] (119) at (-2.5, 3.1) {};
		\node [style=none] (120) at (-4, 3) {};
		\node [style=none] (121) at (-4, 0.475) {};
		\node [style=none] (122) at (-4, -0.875) {};
		\node [style=none] (123) at (-4, -3) {};
		\node [style=none] (124) at (-2.25, -3.125) {};
		\node [style=none] (125) at (-1, -3.25) {};
		\node [style=none] (126) at (-0.05, -3.3) {};
		\node [style=none] (127) at (1.75, -0.35) {};
		\node [style=none] (128) at (1.25, 3.45) {};
		\node [style=none] (129) at (2, 3.525) {};
		\node [style=none] (130) at (2.825, 3.575) {};
		\node [style=none] (131) at (2.775, 0.475) {};
		\node [style=none] (132) at (2.75, -0.775) {};
		\node [style=none] (133) at (2.8, -3.575) {};
		\node [style=none] (134) at (2.05, -3.525) {};
		\node [style=none] (135) at (1.25, -3.45) {};
	\end{pgfonlayer}
	\begin{pgfonlayer}{edgelayer}
		\draw [style=geometric, in=0, out=-90] (3.center) to (5.center);
		\draw [style=geometric, in=90, out=90, looseness=0.75] (12.center) to (13.center);
		\draw [style=dir] (15.center) to (14.center);
		\draw [style=dir] (17.center) to (16.center);
		\draw [in=180, out=-180, looseness=0.50] (49.center) to (48.center);
		\draw [in=180, out=-180, looseness=0.75] (47.center) to (46.center);
		\draw [in=180, out=-180, looseness=0.50] (53.center) to (52.center);
		\draw [in=180, out=-180, looseness=0.50] (51.center) to (50.center);
		\draw [style={open_int}] (57.center) to (58.center);
		\draw [style={open_int}] (59.center) to (60.center);
		\draw [style={open_int}] (83.center) to (82.center);
		\draw [style={open_int}] (84.center) to (85.center);
		\draw [style={open_int}] (61.center) to (86.center);
		\draw [style={open_int}] (87.center) to (88.center);
		\draw [style=geometric, in=180, out=90] (100.center) to (4.center);
		\draw [style=geometric, in=180, out=-90] (100.center) to (5.center);
		\draw [style=geometric, in=180, out=45] (101.center) to (48.center);
		\draw [style=red shade] (100.center)
			 to [in=-150, out=90] (103.center)
			 to [in=-180, out=30, looseness=0.50] (104.center)
			 to [bend left=270, looseness=0.50] (105.center)
			 to [in=45, out=-165] (101.center)
			 to [in=165, out=-45] (106.center)
			 to [in=-165, out=165, looseness=0.50] (107.center)
			 to [in=-30, out=180, looseness=0.75] (102.center)
			 to [in=-90, out=150] cycle;
		\draw [style=red shade] (108.center)
			 to [in=30, out=-90] (111.center)
			 to [in=0, out=-150, looseness=0.50] (112.center)
			 to [bend left=90, looseness=0.50] (113.center)
			 to [in=-150, out=15] (109.center)
			 to [in=-15, out=105] (114.center)
			 to [bend left=90, looseness=0.50] (115.center)
			 to [in=150, out=0, looseness=0.75] (110.center)
			 to [in=90, out=-30] cycle;
		\draw [style=geometric, in=270, out=-105, looseness=0.75] (10.center) to (11.center);
		\draw [style=geometric, in=255, out=-75, looseness=0.75] (92.center) to (7.center);
		\draw [style=geometric, in=120, out=0, looseness=0.75] (48.center) to (116.center);
		\draw [style=red shade] (125.center)
			 to [in=360, out=165, looseness=0.75] (124.center)
			 to (123.center)
			 to [bend left=90, looseness=0.50] (122.center)
			 to [in=-150, out=0, looseness=0.50] (116.center)
			 to [in=-15, out=135, looseness=1.25] (121.center)
			 to [bend left=90, looseness=0.50] (120.center)
			 to [in=-180, out=0] (119.center)
			 to [in=-165, out=0] (118.center)
			 to [in=180, out=-165, looseness=0.25] cycle;
		\draw [style=geometric, in=180, out=-180, looseness=0.25] (55.center) to (126.center);
		\draw [style=geometric, in=90, out=0] (4.center) to (3.center);
		\draw [style=red shade] (135.center)
			 to [in=-180, out=0, looseness=0.75] (134.center)
			 to (133.center)
			 to [bend left=90, looseness=0.50] (132.center)
			 to [in=-45, out=165, looseness=0.75] (127.center)
			 to [in=-165, out=75] (131.center)
			 to [bend left=90, looseness=0.50] (130.center)
			 to [in=0, out=-180] (129.center)
			 to [in=0, out=180, looseness=0.50] (128.center)
			 to [bend right=90, looseness=0.25] cycle;
	\end{pgfonlayer}
\end{tikzpicture}
}}
\end{equation}
\medskip

\nid where the restricted projection $\pi_{|M}\maps M \fib [0,1]$ must be proper and such that $\mathsf{Crit}(\pi_{|M})$ is disjoint from $[0,1]\setminus I$. Note that the red shaded regions of $M$ correspond to the pre-image of $I$ under the restricted projection $\pi_{|M}$, and they contain the critical points of the map --- it is over these regions that the cobordism will undergo meaningful changes in topology as it may contain any finite number of critical points with respect to the projection --- on the other hand, in its complementary regions the property outlined above implies that the embedding must be diffeomorphic to a cylinder of a $(d-1)$-dimensional manifold. Lastly, we note that Calaque--Scheimbauer need their simplicial sets to be Kan complexes so as to make their proof of the Segal condition work, while we do not. We instead rely on the work of Berwick-Evans--Boavida de Brito--Pavlov (see \cite{BBEP24}) to work around this issue, while still maintaining a similar approach as CS to the proof of the Segal condition. 

\emph{Why do we use conically smooth stratified spaces?} 


A central choice in this thesis is to model stratified bordisms using conically smooth stratified spaces, in the sense of Ayala--Francis--Tanaka--Rozenblyum, rather than using more classical theories of stratified spaces, such as Whitney stratified spaces \cite{M12,Tro20}. From our perspective, the point is that conically smooth stratified spaces are singular spaces for which a substantial amount of differential topology survives. They admit a fully developed differential-topological package: one has local models by basics, good notions of smooth map and smooth family, constructible bundles, a robust theory of links, a method called unzipping which allows for functorial resolutions of singularities, tubular neighborhoods, as well as suitably stratified notions of tangent data and tangential structures. In particular, they are especially well tailored to the study of objects in families, which is indispensable for moduli-type constructions like the ones necessitated by modern derived versions of bordism categories. Just as importantly for the purposes of this thesis, this package can be pushed further in directions directly relevant to bordism theory: in Chapter 4 we prove a stratified analogue of the Ehresmann fibration theorem, prove that spaces of conically smooth embeddings into $\mathbb{R}^{\infty}$ are weakly contractible, and establish both weak and depth-$1$ strong Whitney-type embedding results for conically smooth stratified spaces. These results are not ancillary to our later constructions, but rather part of the geometric infrastructure that makes them possible. This is one of the main reasons the conically smooth framework is preferable, for our purposes, to more classical approaches such as Whitney stratifications, where the geometry is often organized through an ambient embedding and is correspondingly less well adapted to intrinsic family-based constructions; see, for instance, the introduction of \cite{AFT1,AFR1}.

Another important point is that conically smooth stratified spaces live within a particularly good categorical context. Rather than just treating stratifications as auxiliary structure on an underlying topological space, they provide a framework in which exit-paths, constructible tangent bundles, and tangential structures fit naturally together. This glue is crucial for the development of stratified bordism categories in the last chapters, where one wants not only to remember how strata meet, but also how singularities vary in families, and how geometric structures behave under pullback. In this respect, the conically smooth formalism is also better suited to our aims than approaches of a more stratified-homotopy-theoretic flavour, such as work of Hughes and others (see \cite{Hu99}), where the emphasis falls instead on homotopy links, path-space constructions, and neighborhood theory. Those approaches are highly natural for questions in stratified homotopy theory, but the present thesis requires a more rigid differential-topological setting, one in which singularities still carry enough geometry to support bordisms, tangential structures, and their families. For this reason, conically smooth stratified spaces provide the most natural ambient category for the constructions undertaken here.

\emph{Why do we restrict to topological manifolds?}
Our construction of stratified bordisms is carried out only for conically smooth stratified spaces whose underlying topological space is a topological manifold of fixed dimension. This is, in part, a conceptual choice. The bordisms relevant for defect topological field theory are meant to model manifold-like bulk geometries together with lower-dimensional strata on which defects live. Requiring the underlying space to remain locally Euclidean ensures that the stratification is encoding controlled singular behaviour inside an ambient topological manifold, rather than replacing the ambient geometry by a more general singular space altogether.

It is also a technical choice. The machinery on which our construction relies — most notably the stratified Ehresmann theorem and the resulting control over families of stratified bordisms — is currently available in this setting, and it is precisely this control that allows the interval Segal machinery to run. We do not view this as a genuine limitation of the theory, only as a first natural range in which the differential-topological and higher-categorical aspects of the construction meet cleanly. Indeed, the formalism developed here appears flexible enough to eventually encompass conically smooth stratified spaces whose underlying space is not a topological manifold.

\section*{Related work}

\nid This thesis lies at the intersection of three bodies of work: the construction of higher bordism categories, the differential topology of conically smooth stratified spaces, and the study of topological field theories with defects. The aim of this section is not to survey these subjects in full, but rather to locate the present construction among the works that are closest either in spirit or in kind.

The foundational backdrop is the standard lineage of higher bordism categories: the embedded topological bordism category of Galatius--Madsen--Tillmann--Weiss \cite{GMTW09}, Lurie's complete $n$-fold Segal-space sketch of the fully extended bordism category \cite{Lur09}, and the detailed Calaque--Scheimbauer construction with tangential structures \cite{CS19}. As explained earlier in this introduction, it is in the Calaque--Scheimbauer framework that the relevant loss-of-fibrancy issue enters: the Segal argument relies on spatial simplicial sets with Kan-complex behaviour, while the smooth singular complexes and moduli spaces involved do not supply this automatically. The construction in this thesis is closest in spirit to this lineage, but uses a new model for $(\oo,1)$-categories which we name interval Segal spaces in contrast to the usual Segal spaces. It makes use of concordance resolutions together with the classifying-space formalism of Berwick-Evans--Boavida de Brito--Pavlov \cite{BBEP24}.

To define interval Segal spaces we introduce a category of intervals $\Int$. This belongs to a broader pattern in which geometric indexing categories replace more rigid combinatorial ones, such as $\Delta$, and one proves that the extra geometric data is homotopically inessential. In stratified factorization homology, Ayala--Francis--Tanaka already relate disk categories to localizations of un-topologized disk categories by inverting isotopy equivalences \cite{AF3}; a close analogue also appears in the work of Karlsson--Scheimbauer--Walde, who prove that the $\oo$-category of disks is a localization of sufficiently fine posets of disks, and relate constructible factorization algebras to the $\oo$-operad of embedded disks \cite{KSW25}. The localization statements for $\Int$ in this thesis should be read in this same spirit: interval data is kept because it is useful for cutting geometric families of bordisms, and then shown not to change the intended homotopy type.

The second source of comparison is the theory of conically smooth stratified spaces developed by Ayala--Francis--Tanaka and Ayala--Francis--Rozenblyum \cite{AFT1,AFR1,AFR2}. Their work provides, not only inspiration, but also the geometric language used throughout: basics, conically smooth maps, constructible bundles, exit-path categories,\footnote{The exit-path category goes back to MacPherson's point of view on constructible sheaves \cite{Mac91}; for published higher-categorical treatments see Lurie \cite{LurieHA} and Treumann \cite{Tre09}.} links, unzipping, and stratified tangential structures. From the perspective of this thesis, conical smoothness serves us not just with a convenient regularity condition but the structure that affords working in families — a primary consideration when extending upwards the ordinary category of stratified bordisms. The theorem of Nocera--Volpe that Whitney stratifications are conically smooth \cite{NV23} further explains why this framework remains close to the classical examples of stratified geometry, while being better adapted to intrinsic moduli constructions. This thesis does not propose a new foundation for stratified spaces. Rather, it uses the conically smooth foundation to construct a symmetric monoidal $(\oo,1)$-category of stratified bordisms and sketch a higher, extended version thereof.

The third source of comparison is the literature on defect topological field theories. Carqueville--Meusburger--Schaumann and Carqueville--Runkel--Schaumann define defect TQFTs as symmetric monoidal functors out of categories of decorated stratified bordisms, and analyze the algebraic structures arising from such theories \cite{CMS20,CRS19}. More recent work, such as the survey of Carqueville--Del Zotto--Runkel \cite{CZR25} and the higher-categorical framework of Müller \cite{Mu25}, places topological defects in a broad field-theoretic context. In comparison, this thesis focuses on the geometric source side. Its goal is not to classify defects algebraically, but to build an $(\oo,1)$-categorical bordism object from which derived defect TQFTs can be defined as symmetric monoidal functors.

Two works are especially close in spirit to the present construction. The first is Will Stewart's thesis on topological domain walls and relative field theories \cite{Ste24}. In that model, the smooth bordism machinery of Calaque--Scheimbauer is adapted to domain walls by replacing ordinary proper submersion conditions with flagged versions thereof. The similarity with our work is that both approaches aim to encode singular or defect-like geometric behavior into bordism constructions. There are, however, two related differences in scope and method. First, Stewart focuses on a restricted and carefully controlled class of domain-wall singularities, whereas the present work builds bordisms directly from conically smooth stratified topological manifolds and allows more general singularity types through the chosen category of basics, with defect labels later encoded through tangential data. Second, Stewart remains close to the flagged version of the Calaque--Scheimbauer model, while this thesis changes the categorical model itself: it uses interval Segal spaces and constructs the spatial levels through concordance resolutions, smooth $\oo$-sheaves, and conically smooth families. Thus the two constructions should be read as closely related but technically distinct approaches to defect-like bordism geometry, rather than as competing treatments of the same model.

The second is the work of Tetik on linked topology \cite{Tet23,Tet22}. Tetik studies depth-$1$ stratified phenomena by replacing a stratified space with a span of ordinary smooth manifolds
\[
M \xleftarrow{\pi} L \xrightarrow{\iota} N,
\]
where the map $\pi$ is a fibre bundle and $\iota$ is a closed embedding. This allows part of the topology of conically smooth stratified spaces, including exit-path quasi-categories and tangential structures, to be expressed using ordinary smooth data. The comparison with conically smooth stratified spaces is already part of Tetik's work: the linked exit-path construction induces a fully faithful functor whose essential image contains the Lurie--MacPherson exit-path quasi-categories of conically smooth spaces over depth-$1$ posets, and the tangential theory of linked smooth manifolds recovers and generalizes the corresponding depth-$1$ conically smooth tangential theory of Ayala--Francis--Rozenblyum. I am not aware of an independent published comparison beyond these results, but this already suggests that linked topology is, at depth $1$, a broader span-based framework in which the conically smooth examples sit naturally. The overlap with this thesis is substantial at the level of motivation: both works are concerned with links, exit paths, tangential structures, and defects. The difference is again one of geometric strategy. Tetik's linked spaces offer a depth-$1$ reduction of stratified geometry to spans of smooth manifolds; this thesis instead keeps conically smooth stratified spaces as the primary geometric objects and constructs an interval Segal space of stratified bordisms from them.

Finally, factorization homology provides a complementary mechanism for producing field theories from local algebraic data. The works of Ayala--Francis, Ayala--Francis--Tanaka, Ayala--Francis--Rozenblyum, and Scheimbauer show how factorization-homological methods organize field theories and observables on manifolds, including stratified manifolds \cite{AF1,AF3,AFR2,Sch14}. In this thesis, factorization homology enters mainly as a future direction: once a suitable stratified bordism category has been constructed, one expects stratified disk-algebras and constructible factorization homology to produce examples of derived defect TQFTs.

I also record several contributions from the unpublished manuscript of Manuel Araújo and Pedro Boavida de Brito, without which the fulfilment of this thesis would not have been possible. The idea of the interval category and its localization as a model for the simplex category underlies Chapter 2. The use of the interval category as a foundation for a variant of complete Segal spaces suited to the $(\oo,1)$-category of bordisms, together with the observation that weak Kan fibrations can be used to transport the proof of \cite[Theorem 5.1]{BBEP24} to the interval Segal condition, enters in Chapter 3. Finally, the notion of the $\cube_n$ category, its localization as a model for the $\bTheta_n$ category, and the use of Segal $\cube_n$-spaces as a foundation for constructing $(\oo,n)$-categories of bordisms enter in Chapter 6.

\section*{Outline}

In Chapter 2 we introduce a model for $(\oo,1)$-categories that complements the existing literature, which we call \emph{interval Segal spaces} — these are simplicial co-presheaves on $\Int$, a subcategory of the poset of opens of $[0,1]$. In particular, we construct a model structure on $\mathsf{coPsh}(\Int)$ which is Quillen equivalent to that of Rezk’s Segal spaces. The motivation here is geometric: most models of $(\oo,1)$-categories are not tailored for the construction of geometric moduli spaces. A rare exception is the recent work of Ayala--Francis--Rozenblyum \cite{AFR1}, where the authors present \emph{striation sheaves} — a highly bespoke class of constructible sheaves on stratified spaces, designed precisely for constructing $(\oo,1)$-categories of a moduli flavour using point-set data. Our construction attempts a similar feat, while being less ambitious and remaining closer to more traditional Segal spaces, thereby facilitating higher-dimensional versions thereof.

In Chapter 3 we present a new construction of the $d$-dimensional bordism $(\oo,1)$-category. This construction is a spiritual cousin of other works, namely those of Galatius--Madsen--Tillman--Weiss \cite{GMTW09}, which later inspired the work of Lurie \cite{Lur09}, and more recently that of Scheimbauer--Calaque \cite{CS19}.

The pioneering work of Galatius--Madsen--Tillman--Weiss offered the first tools to construct an $(\oo,1)$-category of cobordisms. In their paper on the homotopy type of the cobordism category, they construct a topologically enriched category of $d$-dimensional bordisms using classifying spaces. Building on their work, Lurie, in his sketch of a proof of the cobordism hypothesis \cite{Lur09}, put forth a first construction of a model for an $(\oo,n)$-category of cobordisms using $n$-fold complete Segal spaces. This construction contained a small mistake preventing it from satisfying the \emph{essential constancy} condition. In order to address this, in her PhD thesis and in a subsequent paper, Scheimbauer--Calaque propose an alternative definition, which fixes Lurie’s while retaining the use of $n$-fold complete Segal spaces. Despite this almost successful attempt at a \emph{corrigendum}, the work of Scheimbauer--Calaque contains a separate flaw in its claim that the simplicial sets modelling the spatial levels of its $n$-fold Segal spaces are Kan complexes — a claim which, if false, translates into a loss of fibrancy necessary for the proof of the Segal condition.

To remedy this, we make use of the theory of smooth $\oo$-sheaves, i.e. $\oo$-sheaves on the site of smooth manifolds, and their respective classifying spaces — as articulated in the work of Berwick-Evans--Boavida de Brito--Pavlov \cite{BBEP24} — as a conduit for constructing appropriate moduli spaces of smooth manifolds, smooth bordisms, and smooth families of composed bordisms. This gives our approach a distinct flavour and clarifies the minor pathologies identified in previous constructions.

Let us sketch the path we will take in this section. We start by defining a general construction $\widetilde{\bord}_{d,V}(-,-)$ which is a bifunctor on $\man^{\op}\times \Int$ that, when fixing $T\in \Int$, produces a sheaf of sets on $\man^{\op}$ and captures a type of moduli space construction for $d$-dimensional smooth manifolds embedded in $V$ and fibred over smooth manifolds of varying dimension — designed so that, upon taking its concordance resolution and further restricting along the cosimplicial manifold $\bDelta \ra \man$, we obtain a simplicial set $\bord{d,V}$, or space, which we show is the classifying space of $(d-1)$-dimensional manifolds, $d$-dimensional bordisms, and composites of $d$-dimensional bordisms. Lastly, the most dense and technical part of this chapter is spent showing that this construction, a pre-interval space, is an interval Segal space which recovers the $d$-dimensional bordism category $\mathsf{Cob}_d$ as its homotopy 1-category.

In Chapter 4 we finally turn to the theory of \emph{conically smooth stratified spaces}. These were first introduced in the work of Ayala--Francis--Tanaka and further elaborated in the subsequent works of Ayala--Francis--Rozenblyum \cite{AFT1, AFR1, AFR2}. They form the geometric foundations for our stratified higher bordism categories, and their adoption for that role is undoubtedly one of the more novel features of this thesis.

This chapter contains a heavier dose of exposition. We introduce the necessary differential-topological toolkit for our construction of the stratified bordism category, namely: we introduce an $(\oo,1)$-category of conically smooth stratified spaces; the exit-path $(\oo,1)$-category; and the constructible tangent bundle — all of which are necessary for a discussion of stratified tangential structures thereon. We prove new theorems about the contractibility of spaces of embeddings of stratified spaces into Euclidean space, and we prove a stratified analogue of the famed Ehresmann fibration theorem — which is the key to unlocking comparisons with the infrastructure developed for the unstratified case.

In Chapter 5 we construct an interval Segal space, i.e. an $(\oo,1)$-category of conically smooth stratified $d$-dimensional bordisms. We begin by introducing the notion of a conically smooth stratified topological $d$-manifold by stipulating a suitable $(\oo,1)$-category of basics $\Bsc_{d}^{\cS_{d-1}}$, whose links are stratified spheres of appropriate codimensions. Secondly, we produce a pre-interval space $\cS \bord_d$ of stratified topological $d$-dimensional bordisms, again using the concordance resolution for smooth $\oo$-sheaves and a restriction along the cosimplicial smooth manifold $\bDelta \longra \man$, and show that it is an interval Segal space.

Lastly, we endow this interval Segal space with a symmetric monoidal structure agreeing with the expected one coming from disjoint union. After defining this symmetric monoidal interval Segal space, we initiate a discussion on suitable notions of tangential structures for stratified bordisms — in particular, we introduce the notion of \emph{oriented discrete defect datum} by way of discrete right fibrations over $\Bsc_{d}^{\cS_{d-1}}$, as an attempt to make contact with the defect data for stratified bordisms introduced in the work of Carqueville et al. (see \cite{CRS19}) and we prove a comparison result for the stratified oriented case $d=2$.

We conclude Chapter 5 by defining a (derived) \emph{stratified topological quantum field theory} as a symmetric monoidal functor from the $(\oo,1)$-category of stratified bordisms to a symmetric monoidal $(\oo,1)$-category $\cC$.

Chapter 6 is instead prospective in nature. In its first part, following ideas in unpublished work of Manuel Araújo and Pedro Boavida de Brito, we sketch a higher-dimensional generalization of $\Int$, denoted $\cube_d$, with the aim of extending interval Segal techniques to the study of $(\oo,d)$-categories. We conjecture that simplicial co-presheaves on $\cube_d$ satisfying a suitable higher-dimensional Segal condition should present the $(\oo,1)$-category $\Cat_{(\oo,d)}$, and we outline how this formalism could be used to construct a $\cube_d$-space of stratified bordisms, hence providing a candidate model for fully extended stratified topological quantum field theories \`a la Atiyah--Baez--Dolan--Lurie. In the second part of the chapter, we return to the non-extended setting and sketch how stratified factorization homology could provide a mechanism for constructing stratified topological quantum field theories from appropriate stratified disk-algebras, in the spirit of Scheimbauer's factorization-homological approach to topological field theory.
